\subsection{QUOTIENT GROUPS}

PROPOSITION 16. Let G be a topological group and let H be a normal subgroup of G. Then the quotient by H of the topology of G is compatible with the group structure of G/H.

If \(x \to \dot{x}\) is the canonical mapping of G onto G/H, then we have to show that \((\dot{x}, \dot{y}) \to \dot{x}\dot{y}^{-1}\) is a continuous mapping of (G/H) \(\times\) (G/H) into G/H. Since the equivalence relation \(x^{-1}y \in H\) is open (no. 4, Lemma 2), it is enough to show that \((x,y) \to \dot{x}\dot{y}^{-1}\) is a continuous mapping of G \(\times\) G into G/H (Chapter I, § 5, no. 3, Corollary to Proposition 8, and § 3, no. 4, Proposition 6). But \((x,y) \to \dot{x}\dot{y}^{-1}\) is the composition of the continuous mappings \(x \to \dot{x}\) and \((x,y) \to xy^{-1}\) , hence is continuous.

Whenever in the sequel we consider a quotient group G/H of a topological group G as a topological group, it is always to be understood that the topology of G/H is the quotient by H of the topology of G, unless the contrary is expressly stated.

PROPOSITION 17. Let \(\varphi\) be the canonical mapping of a topological group \(\mathbf{G}\) onto a quotient group \(\mathrm{G / H}\) . If \(\mathfrak{B}\) is a fundamental system of neighbourhoods of \(e\) in \(\mathbf{G}\) , then \(\varphi(\mathfrak{B})\) is a fundamental system of neighbourhoods of the identity element \(\varphi(e)\) of \(\mathrm{G / H}\) .

This is a particular case of Proposition 5 of Chapter I, § 5, no. 3.

Propositions 13 and 14 give in particular, for quotient groups:

PROPOSITION 18. Let G be a topological group and let H be a normal subgroup of G.

\begin{enumerate}
\def\labelenumi{\alph{enumi})}
\item
  The quotient group G/H is Hausdorff if and only if H is closed in G.
\item
  The quotient group G/H is discrete if and only if H is open in G.
\end{enumerate}

If G is a topological group and N is the closure of \(\{e\}\) in G, then N is a closed normal subgroup of G (no. 1, Proposition 1), hence G/N is Hausdorff; G/N is called the Hausdorff group associated with G.

PROPOSITION 19. If H is a discrete normal subgroup of a topological group G, then G/H is locally isomorphic to G.

Let \(\mathbf{V}\) be a neighbourhood of \(e\) in \(\mathbf{G}\) which does not contain any point of \(\mathbf{H}\) other than \(e\) , and let \(\mathbf{W}\) be a symmetric open neighbourhood of \(e\) in \(\mathbf{G}\) such that \(\mathbf{W}^2 \subset \mathbf{V}\) . Then the restriction to \(\mathbf{W}\) of the canonical mapping \(\varphi\) of \(\mathbf{G}\) onto \(\mathbf{G} / \mathbf{H}\) is injective; for if \(x, y \in \mathbf{W}\) and \(\varphi(x) = \varphi(y)\) , then \(x^{-1}y \in \mathbf{W}^2 \subset \mathbf{V}\) and \(x^{-1}y \in \mathbf{H}\) , so that \(x = y\) . By Proposition 17 it follows that the restriction of \(\varphi\) to \(\mathbf{W}\) is a homeomorphism of \(\mathbf{W}\) onto \(\varphi(\mathbf{W})\) ; since moreover \(\varphi(xy) = \varphi(x)\varphi(y)\) for all \(x, y \in \mathbf{W}\) , we conclude that \(\mathbf{G}\) and \(\mathbf{G} / \mathbf{H}\) are locally isomorphic (§ 1, no. 3, Proposition 3).

